%==============================================================================
%== template for LATEX poster =================================================
%==============================================================================
%
%--A0 beamer slide-------------------------------------------------------------
\documentclass[final]{beamer}
\usetheme{Warsaw}
\usepackage[utf8]{inputenc}
\usepackage[T2A]{fontenc}
 \usepackage[english,russian]{babel}
 \usepackage{amsmath,amssymb}
 \usepackage{amscd}
\usepackage{amsthm}
\usepackage[orientation=portrait,size=a1,
            scale=1.25         % font scale factor
           ]{beamerposter}

\geometry{
  hmargin=2.5cm,
}



\linespread{1.15}

\usetheme{sharelatex}

\title[] % Если эта запись не нужна, удалите всё внутри квадратных скобок.
%
{\bf \huge % Poster title
Название. Если слишком длинное,
\texorpdfstring{\vspace{1cm}}{}
то на две строки
}

\author{ % Authors
Иванов Иван Иванович
}

 % General University
 \institute{
ФКН НИУ ВШЭ, Москва
}

\date{\today}

\begin{document}

\begin{frame}[t]
%==============================================================================
\begin{multicols}{2} %Если не зхватает места, можно изменить 2 на 3, будет 3 колонки. 
%==============================================================================
%==The poster content==========================================================
%==============================================================================

\section{Основные понятия}


Здесь введены основные определения. Также рассказаны базовые факты про них. Без лишних подробностей, но так, чтобы математик, не являющийся специалистом в узкой теме доклада, мог полностью понять дальнейшее.


Безусловно, можно писать формулы 

\begin{equation}
H = \sum_{i=1}^{N} h_{D}(i) + \sum_{j>i=1}^{N} C_{ij}
\end{equation}

или ссылаться на литературу, такую как~\cite{ref1} или~\cite{ref1,ref2}...

\vspace{0.8cm}
\structure{\underline{Определение.}} Особо важным {\color{color2}объектам} стоит дать формальное определение.
\vspace{0.8cm}

Стиль оформления стоит стараться сохранить. 




\section{Историческая справка}

В работах~\cite{ref1} и~\cite{ref2} была доказана следующая теорема.

\vspace{0.8cm}
\structure{\underline{Теорема 1.}} Для $n=1$ и $2$ выполнены прекрасные свойства.
\vspace{0.8cm}

В работе~\cite{ref2} высказана гипотеза.  

\vspace{0.8cm}
\structure{\underline{Гипотеза.}} Прекрасные свойства выполнены для любого $n$.
\vspace{0.8cm}

В нашей работе доказана гипотеза для чётных $n$ вида $8k+7$. 

\section{Основные результаты}

Основным результатом работы является доказательство следующих утверждений. 

\vspace{0.8cm}
\structure{\underline{Предложение.}} Если число делится на 4, то оно возможно чётное.
\vspace{0.8cm}

\structure{\underline{Теорема 2.}} Если число не делится на 4, то оно возможно нечётное.
\vspace{0.8cm}

Несмотря на простоту формулировок, эти утверждения имеют очень технически сложные и довольно длинные доказательства. В них используется теория эллиптических кривых, фракталы, а также устный счёт в пределах от 1 до 10. Если использовать теорему Гёделя о неполноте и вероятностные методы, то можно использовать устный счёт только в пределах 8, что по силам среднестатистическому филдсовскому медалисту.

\subsection{Идеи доказательств}


Как видно из таблицы 

\vskip1ex
\begin{table}
\centering
\caption{Буквы и цифры.}
\begin{tabular}{ccccc}
\hline\hline
1 & 2 & 3 & 4 & 5\\
\hline
aaa & bbb & ccc & ddd & eee\\
aaaa & bbbb & cccc & dddd & eeee\\
aaaaa & bbbbb & ccccc & ddddd & eeeee\\
aaaaaa & bbbbbb & cccccc & dddddd & eeeeee\\
1.000 & 2.000 & 3.000 & 4.000 & 5.000\\
\hline\hline
\end{tabular}
\end{table}
\vskip1ex


теорема~1 выполняется для всех чётных чисел вида $8k+9$. Далее применяем логарифмирование с последующим взятием экспоненты. Поскольку процесс забавный, итерируем его много раз. Как видим, ничего не происходит. Отсюда и следуют наши утверждения.  


\section{Примеры}

Так как чётных чисел вида $8k+7$ не бывает, то и примеров у нашей теоремы нет. Зато мы можем предъявить отдельно чётное число и отдельно число вида $8k+7$. В принципе, эти 2 примера в достаточной степени иллюстрируют теорему и в лучших примерах она не нуждается. 


Далее следует длинный текст. 
Далее следует длинный текст. 
Далее следует длинный текст. 
Далее следует длинный текст. 
Далее следует длинный текст. 
Далее следует длинный текст. 
Далее следует длинный текст. 
Далее следует длинный текст. 
Далее следует длинный текст. 
Далее следует длинный текст.
Далее следует длинный текст. 
Далее следует длинный текст. 
Далее следует длинный текст. 
Далее следует длинный текст. 
Далее следует длинный текст.
Далее следует длинный текст. 
Далее следует длинный текст. 
Далее следует длинный текст. 
Далее следует длинный текст. 
Далее следует длинный текст.
Далее следует длинный текст. 
Далее следует длинный текст. 
Далее следует длинный текст. 
Далее следует длинный текст. 
Далее следует длинный текст.




\section{Выводы и дальнейшие направления исследований}

Наш результат крайне важен для развития науки. Иначе как бы нам удалось сделать по нему постер на этой конференции. 

\vskip1ex
\begin{figure}
\centering
\includegraphics[width=0.9\columnwidth]{SCFCS-Poster-Template-2025/fcs-logo-wide.png}
\caption{Логотип ФКН.}
\end{figure}
\vskip1ex

В дальнейшем мы собираемся доказать нашу теорему для чётных $n$ вида $8k+5$, а также для нечётных $n$ вида $8k+2$. И тут внезапно очень захотелось подчеркнуть, что центральным в этой работе было понятие {\color{color4}чётности}.

Далее следует длинный текст. 
Далее следует длинный текст. 
Далее следует длинный текст. 
Далее следует длинный текст. 
Далее следует длинный текст. 
Далее следует длинный текст. 
Далее следует длинный текст. 
Далее следует длинный текст. 
Далее следует длинный текст. 
Далее следует длинный текст.
Далее следует длинный текст. 
Далее следует длинный текст. 
Далее следует длинный текст. 
Далее следует длинный текст. 
Далее следует длинный текст.
Далее следует длинный текст. 
Далее следует длинный текст. 
Далее следует длинный текст. 
Далее следует длинный текст. 
Далее следует длинный текст.
Далее следует длинный текст. 
Далее следует длинный текст. 
Далее следует длинный текст. 
Далее следует длинный текст. 
Далее следует длинный текст.
Далее следует длинный текст. 
Далее следует длинный текст. 
Далее следует длинный текст. 
Далее следует длинный текст. 
Далее следует длинный текст.
Далее следует длинный текст. 
Далее следует длинный текст. 
Далее следует длинный текст. 
Далее следует длинный текст. 
Далее следует длинный текст.
Далее следует длинный текст. 
 




%==============================================================================
%==End of content==============================================================
%==============================================================================

%--References------------------------------------------------------------------

\subsection{Список литературы}

\begin{thebibliography}{99}

\bibitem{ref1} C.P. Ramanujam. ''A note on automorphism group of algebraic variety''. Math. Ann. 156 (1964), 25-33.

\bibitem{ref2} Э. Б. Винберг, ''Группа Вейля градуированной алгебры Ли'', Изв. АН СССР. Сер. матем., 40:3 (1976), 488–526.

\end{thebibliography}
%--End of references-----------------------------------------------------------

\end{multicols}

%==============================================================================
\end{frame}
\end{document}
